\documentclass[12pt]{article}
\usepackage[utf8]{inputenc}
\usepackage[T1]{fontenc}
\usepackage{hyperref}
\title{Tzotzil, Huixt\'an Language Resources for K-12 Curriculum Development}
\author{Compiled by Claude Sonnet 5 (claude-sonnet-5) for Andr\'as Kornai}
\date{}
\begin{document}
\maketitle

\section{Introduction}
This report covers Tzotzil of Huixt\'an (ISO 639-3 tzu, Glottocode tzot1261). Little exists for this variety, and most catalog entries file it under the wider Tzotzil code (tzo).

\section{Primary Educational Resources}
\begin{itemize}
\item A 1970 booklet of Tzotzil tales (16 pages) in the Huixt\'an variety, catalogued in the Colegio de M\'exico repository. The catalog notes that its language is no longer current. \url{https://repositorio.colmex.mx/concern/books/th83m1062?locale=en}
\item A 1973 bilingual Spanish--Huixt\'an sentence book, published with SIL, in the same repository. \url{https://repositorio.colmex.mx/concern/books/dz010r91g?locale=en}
\end{itemize}

\section{Linguistic and Reference Materials}
\begin{itemize}
\item The ASJP database page for Tzotzil Huixt\'an cites Cowan's 1969 \emph{Tzotzil Grammar} (SIL) as a source. \url{https://asjp.clld.org/languages/TZOTZIL_HUIXTAN}
\end{itemize}

\section{Religious and Community Resources}
Joshua Project lists Scripture Earth gospel resource links for Huixt\'an and related varieties. I did not see the Scripture Earth page itself. \url{https://joshuaproject.net/index.php/languages/tzo}

\section{Digital and Audio Resources}
\begin{itemize}
\item Global Recordings Network audio in Huixt\'an Tzotzil: evangelism and Bible teaching recordings, an MP3 story, and a slideshow. \url{https://globalrecordings.net/en/language/1204}
\end{itemize}

\section{Curriculum Development Context}
I found no modern dictionary, school programme or curriculum for this variety. The available materials are mostly missionary and older print items, and would need considerable new development work.

\end{document}